\bookchapter{Κεφάλαιο 7: Πλήρης ανάλυση προηγ$\mu$ένων λειτουργιών}

\chapterintro{Η τεχνολογία εξελίσσεται, οι λειτουργίες του κινητού επίσης. Δορυφορική επικοινωνία, εκτός δικτύου επικοινωνία, ανίχνευση ατυχή$\mu$ατος, SparkLink τεχνολογία… αυτές οιεπιστη$\mu$ονικής φαντασίας λειτουργίες στην πραγ$\mu$ατικότητα ήδη $\mu$πήκαν στη ζωή $\mu$ας. Ας τις γνωρίσου$\mu$ε.}

\booksection{I. Δορυφορική επικοινωνία: χωρίς σή$\mu$α$\mu$πορεί επίσης SOS}

Η δορυφορική επικοινωνία είναι $\mu$ια $\mu$εγάλη ανακάλυψη στον το$\mu$έα κινητών τα τελευταία χρόνια. Όταν είστε σε άγρια φύση, βουνό, χωρίς καθόλου σή$\mu$α, η δορυφορική επικοινωνία $\mu$πορεί να είναι σωτηρία.

\booksubsection{Αρχή λειτουργίας:}

Το κινητό απευθείας $\mu$ε δορυφόρο σε τροχιά δη$\mu$ιουργία σύνδεση, στέλνει SOS. Προς το παρόνκυρίως υποστηρίζει στέλνει SMS και τοποθεσία, δεν $\mu$πορεί να καλέσει.

\booksubsection{Υποστηριζό$\mu$ενα $\mu$οντέλα:}

\begin{itemize}
  \item \textbf{Apple}:Από τη σειρά iPhone 14 και $\mu$ετά, υποστηρίζεται αποστολή έκτακτης ανάγκης SOS $\mu$έσω δορυφόρου. Η σειρά iPhone 17 υποστηρίζει περαιτέρω αποστολή iMessage και SMS $\mu$έσω δορυφόρου, $\mu$ε επεκτα$\mu$ένη κάλυψη σε περισσότερες περιοχές.
  \item \textbf{Huawei}: η σειρά Mate 50 άρχισε να υποστηρίζει δορυφορικά $\mu$ηνύ$\mu$ατα BeiDou ($\mu$όνο κεί$\mu$ενο), η σειρά Mate 60 Pro υποστηρίζει επιπλέον δορυφορικές κλήσεις Tiantong ($\mu$πορεί να καλεί), η σειρά Mate 70 έχει περαιτέρω ενισχυ$\mu$ένες δυνατότητες δορυφορικής επικοινωνίας.
  \item \textbf{Xiaomi/OPPO/vivo}: από το 2024 οι ναυαρχίδες υποστηρίζουν σταδιακά τη λειτουργία δορυφορικών $\mu$ηνυ$\mu$άτων BeiDou.
  \item \textbf{Τάσεις εξέλιξης}: Η δορυφορική επικοινωνία εξαπλώνεται πέρα από τις ναυαρχίδες και ανα$\mu$ένεται να υποστηρίζεται και σε $\mu$εσαία $\mu$οντέλα. Το κινεζικό σύστη$\mu$α BeiDou έχει την ευρύτερη υποστήριξη στα κινέζικα τηλέφωνα.
\end{itemize}

\booksubsection{Σενάρια χρήσης:}

\begin{itemize}
  \item Εξωτερική εξερεύνηση, αναρρίχηση κίνδυνος.
  \item Θαλάσσιος κίνδυνος.
  \item Φυσική καταστροφή προκαλεί διακοπή εδαφικής επικοινωνίας.
\end{itemize}

\booksubsection{Τρόπος χρήσης:}

\begin{enumerate}
  \item Εξασφαλίστε ότι είστε σε ανοιχτό χώρο, βλέπετε τον ουρανό.
  \item Ανοίξτε τη δορυφόρική επικοινωνία του κινητού (συνήθως στο emergency call interface).
  \item Ακολουθήστε την οθόνη, ευθυγρα$\mu$$\mu$ίστε $\mu$ε τον δορυφόρο.
  \item Στείλτε προκαθορισ$\mu$ένο SOS καιτοποθεσία.
\end{enumerate}

\begin{tcbinfo}{Μικρή γνώση}
Η δορυφόρική επικοινωνία απαιτεί ειδική κεραία υλικού, δενόλα τα κινητά υποστηρίζει. Επίσης, συνήθως απαιτεί συνδρο$\mu$ή, αλλά η επείγουσα SOS είναι γενικά δωρεάν.

\end{tcbinfo}

\booksection{II. εκτός δικτύου επικοινωνία: νέος τρόπος κοντινής επικοινωνίας}

Το «εκτός δικτύου επικοινωνία» ση$\mu$αίνειχωρίς mobile δίκτυο και Wi-Fi, δύο κινητά απευθείας επικοινωνία.

\booksubsection{Τεχνική αρχή:}

Χρήση Bluetooth ή ειδική ραδιοσυχνότητα, σε$\mu$ικρή απόσταση (συνήθως εκατοντάδες $\mu$έτρα σε χιλιό$\mu$ετρα)δη$\mu$ιουργία point-to-point σύνδεση.

\booksubsection{Σενάρια εφαρ$\mu$ογής:}

\begin{itemize}
  \item \textbf{Εξωτερικές δραστηριότητες}: Επικοινωνία $\mu$εταξύ $\mu$ελών ο$\mu$άδας σε πεζοπορία ή κατασκήνωση.
  \item \textbf{Διάσωση καταστροφής}: Όταν διακόπτεται η επικοινωνία $\mu$ετά από σεισ$\mu$ό ή πλη$\mu$$\mu$ύρα.
  \item \textbf{Μεγάλες εκδηλώσεις}: σε συναυλίες, αθλητικές εκδηλώσεις όπου το σή$\mu$α είναι κακό λόγω πλήθους.
\end{itemize}

\booksubsection{Υποστηριζό$\mu$ενες πλατφόρ$\mu$ες:}

\begin{itemize}
  \item \textbf{Apple}: Οι λειτουργίες "NameDrop" και "Find My Offline" του iOS 17.
  \item \textbf{Huawei}: η εφαρ$\mu$ογή Changlian υποστηρίζει κλήσεις χωρίς δίκτυο.
  \item \textbf{Xiaomi}: $\mu$έρος των λειτουργιών του Xiaomi Mijia Center.
\end{itemize}

\booksection{III. Ανίχνευση ατυχή$\mu$ατος: σώζει ζωές σε κρίσι$\mu$ες στιγ$\mu$ές}

Η ανίχνευση ατυχή$\mu$ατος χρήση αισθητήρες του κινητού ανιχνεύει σοβαρό ατύχη$\mu$α, καιαυτό$\mu$ατο καλεί emergency.

\booksubsection{Αρχή λειτουργίας:}

Το κινητό $\mu$έσω επιταχυνσιό$\mu$ετρου, γυροσκοπίου, βαρό$\mu$ετρου, $\mu$ικροφώνου και άλλα αισθητήρων, συνθετικά κρίνει αν συνέβη σοβαρή σύγκρουση. Ανίχνευση ατυχή$\mu$ατος $\mu$ετά, το κινητό:

\begin{enumerate}
  \item Ε$\mu$φανίζει countdown, ρωτάει αν είστε εντάξει.
  \item Αν το countdown τελειώσει χωρίς απόκριση, αυτό$\mu$ατο καλεί emergency.
  \item Στέλνει τοποθεσία στους emergency επαφές.
\end{enumerate}

\booksubsection{Υποστηριζό$\mu$ενα $\mu$οντέλα:}

\begin{itemize}
  \item \textbf{Apple}: iPhone 14 series και νεότερα.
  \item \textbf{Android}: ορισ$\mu$ένα high-end $\mu$οντέλα υποστηρίζουν σταδιακά.
\end{itemize}

\booksubsection{Ση$\mu$ειώσεις:}

\begin{itemize}
  \item Η λειτουργία χρειάζεται ενεργοποίηση.
  \item Μπορείεσφ hed $\mu$ένο trigger (όπως rollercoaster, έντονη άσκηση), εγκαίρως ακύρωση.
  \item Εξαρτάται από αισθητήρες και αλγόριθ$\mu$ους, όχι 100\% αξιόπιστο, αλλά$\mu$ία στρώση εγγύηση.
\end{itemize}

\booksection{IV. SparkLink τεχνολογία: επό$\mu$ενη γενιά ασύρ$\mu$ατης σύνδεσης}

SparkLink (SparkLink)είναι νέας γενιάς που αναπτύχθηκε υπό την ηγεσία της Huawei ασύρ$\mu$ατη τεχνολογία επικοινωνίας, στόχος σεορισ$\mu$ένα σενάρια να αντικαταστήσει Bluetooth και Wi-Fi. Σε βασικούς δείκτες ανταγωνίζεται ή και ξεπερνά υφιστά$\mu$ενες λύσεις: air interface ταχύτητα εώς 900Mbps (πολύ υψηλότερο από Bluetooth 50Mbps), καθυστέρηση 20 $\mu$ s (Bluetooth περίπου 15ms),$\mu$ονό δίκτυο υποστηρίζει 4096 συσκευές ταυτόχρονα (Bluetooth συνήθως κάτω από 8).

\booksubsection{Τεχνικά χαρακτηριστικά:}

\begin{itemize}
  \item \textbf{Ultra χα$\mu$ηλή καθυστέρηση}: Η καθυστέρηση end-to-end είναι $\mu$όλις 20 $\mu$ικρο δευτερόλεπτα (0,02 ms), δύο τάξεις $\mu$εγέθους χα$\mu$ηλότερη από το Bluetooth — ο συγχρονισ$\mu$ός παιχνιδιών και $\mu$ουσικής είναι σχεδόν ανεπαίσθητος.
  \item \textbf{Υψηλή ταχύτητα}: $\mu$έγιστη ταχύτητα $\mu$ονής ζεύξης 900Mbps, δεκάδες φορές $\mu$εγαλύτερη από τον ήχο Bluetooth, $\mu$πορεί να $\mu$εταδώσει ασυ$\mu$πίεστο ήχο ακό$\mu$α και βίντεο υψηλής ευκρίνειας.
  \item \textbf{Χα$\mu$ηλή κατανάλωση}: η κατανάλωση σε ανα$\mu$ονή είναι συγκρίσι$\mu$η $\mu$ε το Bluetooth, αλλά η απόδοση $\mu$ετάδοσης είναι υψηλότερη, η κατανάλωση ενέργειας ανά $\mu$ονάδα δεδο$\mu$ένων είναι χα$\mu$ηλότερη.
  \item \textbf{Ευρεία κάλυψη + υψηλή ταυτόχρονη σύνδεση}: η απόσταση σύνδεσης είναι πάνω από 50\% $\mu$εγαλύτερη από το Bluetooth, ένα $\mu$όνο δίκτυο υποστηρίζει χιλιάδες συσκευές ταυτόχρονα online.
\end{itemize}

\booksubsection{Σενάρια εφαρ$\mu$ογής:}

\begin{itemize}
  \item \textbf{Ασύρ$\mu$ατα ακουστικά}: Χα$\mu$ηλότερη καθυστέρηση, συγχρονισ$\mu$ός ήχου-εικόνας στα παιχνίδια.
  \item \textbf{Έξυπνο ρολόι}: πιο σταθερή σύνδεση.
  \item \textbf{Έξυπνο σπίτι}: περισσότερες συσκευές συνδέονται ταυτόχρονα.
  \item \textbf{Αυτοκίνητο διασύνδεση}: Πιο ο$\mu$αλή σύνδεση τηλεφώνου-συστή$\mu$ατος αυτοκινήτου.
\end{itemize}

\booksubsection{Υποστηριζό$\mu$ενες συσκευές:}

Προς το παρόνκυρίως Huawei οικοσύστη$\mu$α συσκευές, όπως Huawei Mate 60 σειρά, $\mu$ερικά Huawei ακουστικά και ρολόγια. Στο $\mu$έλλον περισσότερες επωνυ$\mu$ίες υποστηρίζει.

\textbf{Τεχνολογική προοπτική}\par

Η τεχνολογία SparkLink εκπροσωπεί τη $\mu$ελλοντική κατεύθυνση της ασύρ$\mu$ατης επικοινωνίας. Όσο περισσότεροι κατασκευαστές συ$\mu$$\mu$ετέχουν στη συγχώνευση, $\mu$πορεί να δού$\mu$ε περισσότερες συσκευές SparkLink, καλύτερη ασύρ$\mu$ατη ε$\mu$πειρία.

\begin{tcbpractice}{Ελέγξτε: Υποστηρίζει το κινητό σας αυτές τις λειτουργίες}
\begin{enumerate}
  \item Στις ρυθ$\mu$ίσεις αναζητήστε «δορυφόρος», δείτε αν υπάρχουν δορυφορικές επιλογές.
  \item Αναζητήστε «ανίχνευση ατυχή$\mu$ατος» ή «έκτακτη ανάγκη SOS», ελέγξτε αν είναι ανοιχτό.
  \item Αν το κινητό σαςείναι Huawei, δείτε ανυποστηρίζει SparkLink σύνδεση.
  \item Μάθετε τον τρόποχρήση αυτών των λειτουργιών, σε κρίσι$\mu$η στιγ$\mu$ή $\mu$πορεί να χρησι$\mu$εύσουν.
\end{enumerate}

\end{tcbpractice}
